\subsection{既约代数与平展代数}

定义 2. —— 设 A 为一个交换环；若 A 中每个幂零元（I，第 98 页）均为零，则称 A 为既约的。

如果 A 是一个域、一个整环或既约环的乘积，那么它是一个既约环。要使交换环 A 是既约的，必要且充分的是：\(a^{2} \neq 0\) （对 A 中每个 \(a \neq 0\)）：因为由此通过对 n 归纳我们可得 \(a^{2^{n}} \neq 0\) ，因此 \(a^{n} \neq 0\) （对 A 中任意 \(a \neq 0\)）。

一个代数称为既约的，如果其底环是既约的。

命题5。——设 A 是 K 上有限次数的交换代数。要使 A 是既约的，必要且充分的是：存在 K 上的有限次扩张 \(L_{1}, \ldots, L_{n}\) ，使得 A K-同构于 \(L_{1} \times \cdots \times L_{n}\) 。

所述条件显然充分。

反之，假设 A 是既约的；对 A 的次数进行归纳论证，我们看到只需证明：若 A 不是域，则存在两个非零代数 \(A_{1}\) 与 \(A_{2}\) 使得 A 同构于 \(A_{1} \times A_{2}\) ，或者等价地说，A 中存在一个不同于 0 和 1 的幂等元。

从现在起假设 A 是既约的且不是域。在 A 中不同于 0 和 A 的理想里，设 a 是作为向量 K-空间维数最小的理想。对于 a 中任意 \(x \neq 0\) 我们有 \(x^{2} \neq 0\) ，因为 A 是既约的，因此 \(a^{2} \neq \{0\}\) 。我们有 \(a^{2} \subset a\) ，并且由 a 的最小性，我们得到 \(a^{2} = a\) 。根据引理 4，存在幂等元 e 使得 a = Ae，并且我们有 \(e \neq 0\) , \(e \neq 1\) ，因为 a 不同于 0 和 A。

定理 4. --- 设 A 是 K 上有限次数的交换代数。则以下断言等价：

\begin{enumerate}
\def\labelenumi{\alph{enumi})}
\item
  该代数 \(\mathbf{A}\) 是平展的。
\item
  对于 K 的每个扩张 L，环 L ⊗K A 是既约的。
\item
  存在一个完全扩张域 \(\mathbf{P}\) （\(\mathbf{K}\) 的），使得环 \(\mathbf{P} \otimes_{\mathbf{K}} \mathbf{A}\) 是既约的。
\end{enumerate}

* d) 存在可分代数扩张 \(\mathbf{L}_1, \ldots, \mathbf{L}_n\) （\(\mathbf{K}\) 的），使得 \(\mathbf{A}\) 同构于 \(\mathbf{L}_1 \times \cdots \times \mathbf{L}_n\) 。*

特别地，每个平展代数都是既约的。

\begin{enumerate}
\def\labelenumi{\Alph{enumi})}
\tightlist
\item
  我们首先证明 \(a\) , \(b\) 与 \(c\) 的等价性。
\end{enumerate}

假设 A 是平展的，并设 L 为 K 的一个扩域。设 \(\Omega\) 为 L 的一个代数闭扩域（V，第 23 页，定理 2）。则 \(L \otimes_{K} A\) 同构于 \(\Omega \otimes_{K} A\) 的一个子环，且后者由命题 2（V，第 30 页）同构于一个环 \(\Omega^{n}\) 。因此环 \(L \otimes_{K} A\) 是既约的。

我们从而证明了 \(a)\) 蕴含 \(b)\) ，而 \(c)\) 是 \(b)\) 的一个特殊情形。现在假设 \(c)\) 成立。要使 K-代数 A 是平展的，必要且充分的是 P-代数 \(\mathbf{A}_{(P)}\) 是平展的（V，第 32 页，推论 2）。因此由以下引理，代数 A 是平展的：

引理 5. —— 设 B 为完全域 P 上的有限次既约代数；则 B 是平展的。

由命题 5，存在扩域 \(L_{1}, \ldots, L_{n}\) （P 的），使得 B 同构于代数 \(L_{1} \times \cdots \times L_{n}\) 。由于有限个平展代数的乘积是平展的（V，第31页，推论），因此只需考察 B 是 P 的扩张的情形。根据定理3（V，第33页），只需证明在 \(\Omega_{\mathrm{P}}(\mathrm{B})\) 中对所有 \(x \in B\) 有 dx = 0。

设 \(x \in B\) ；由于 \(B\) 在 \(K\) 上是有限次的，\(x\) 在 \(K\) 上是代数的（V，第18页，命题2）。设 \(f\) 为 \(x\) 的极小多项式，并令 \(f'\) 为 \(f\) 的导数。多项式 \(f\) 是非常数的。假设 \(f' = 0\) ；由 \(V\) 第9页推论，域 \(P\) 的特征为 \(p \neq 0\) ，且我们有 \(f \in P[X^p]\) ；由于 \(P\) 是完全的，我们有 \(P[X^p] = P[X]^p\) ，但不可约多项式 \(f\) 不能属于 \(P[X]^p\) 。

因此我们有 \(f' \neq 0\) ，且由于 \(f'\) 的次数严格小于 f 的次数，我们有 \(f'(x) \neq 0\) 。现在由 \(f(x) = 0\) 我们推出 \(f'(x) \cdot dx = 0\) 在 \(\Omega_{\mathrm{P}}(\mathrm{B})\) ，从而 dx = 0，正如我们所要证明的。

* B) 假设 A 是平展的；由 a) 与 b) 的等价性，代数 A 是既约的，因此存在扩张 \(L_{1}, \ldots, L_{n}\) （K 的），使得 A 同构于 \(L_{1} \times \cdots \times L_{n}\) （命题 5）。由于 A 是平展的，每个扩张 \(L_{i}\) 都是平展代数（V，第～31页，推论），因此按定义是 K 的可分代数扩张。

蕴含关系 \(d) \Rightarrow a)\) 由 V 第～31页推论得出。*

推论。--- 假设 K 的特征为 \(p \neq 0\) 。要使 A 是平展的，必要且充分的是 \(A = K[A^{p}]\) 。此时，对 A 在 K 上的每个基 \((a_{i})_{i \in I}\) ，族 \((a_{i}^{p})_{i \in I}\) 是 A 在 K 上的一个基。

我们选取 \(\Omega\) （K 的一个代数闭包）。给定两个 K-同态 \(u\) 与 \(v\) （从 A 到 \(\Omega\)），如果 \(u\) 与 \(v\) 在 \(\mathbf{K}[\mathbf{A}^p]\) 上的限制相同，则我们有

\[
u (x) ^ {p} = u \left(x ^ {p}\right) = v \left(x ^ {p}\right) = v (x) ^ {p}, \quad \text { whence } \quad u (x) = v (x)
\]

对所有 \(x \in \mathbf{A}\) 。于是我们有不等式 \([\mathbf{A} : \mathbf{K}]_s \leqslant [\mathbf{K}[\mathbf{A}^p] : \mathbf{K}]_s\) ；若 \(\mathbf{A}\) 是平展的，则我们有

\[
[ \mathrm{A}: \mathrm{K} ] = [ \mathrm{A}: \mathrm{K} ] _ {s} \leqslant [ \mathrm{K} [ \mathrm{A} ^ {p} ]: \mathrm{K} ] _ {s} \leqslant [ \mathrm{K} [ \mathrm{A} ^ {p} ]: \mathrm{K} ],
\]

由此 \(\mathbf{A} = \mathbf{K}[\mathbf{A}^p]\) 。

反之，假设我们有 \(\mathbf{A} = \mathbf{K}[\mathbf{A}^p]\) ；设 \((a_i)_{i\in I}\) 是 \(\mathbf{A}\) 在 \(\mathbf{K}\) 上的一个基。根据 \(\mathbf{V}\) 第4页命题2， \(b\) ），族 \((a_i^p)_{i\in I}\) 生成向量 \(\mathbf{K}\) -空间 \(\mathbf{K}[\mathbf{A}^p]\) ，且由于 \(\mathbf{A} = \mathbf{K}[\mathbf{A}^p]\) 的维数有限并等于 \(\mathbf{I}\) 的基数，族 \((a_i^p)_{i\in I}\) 是 \(\mathbf{A}\) 在 \(\mathbf{K}\) 上的一个基。设 \(u\) 为 \(\Omega \otimes_{\mathbf{K}} \mathbf{A}\) 的一个元素，使得 \(u^2 = 0\) ，从而 \(u^p = 0\) ；由于 \((a_i)_{i\in I}\) 是 \(\mathbf{A}\) 在 \(\mathbf{K}\) 上的一个基，存在一个族 \((\lambda_i)_{i\in I}\) （其元素属于 \(\Omega\) ），使得 \(u = \sum_{i\in I} \lambda_i \otimes a_i\) ，从而 \(u^p = \sum_{i\in I} \lambda_i^p \otimes a_i^p\) 。由于 \((a_i^p)_{i\in I}\) 是 \(\mathbf{A}\) 在 \(\mathbf{K}\) 上的一个基，且 \(u^p = 0\) ，我们必须有 \(\lambda_i^p = 0\) ，从而 \(\lambda_i = 0\) （对所有 \(i\in I\)），因此 \(u = 0\) 。这表明环 \(\Omega \otimes_{\mathbf{K}} \mathbf{A}\) 是既约的；由于 \(\Omega\) 是完全的，代数 \(\mathbf{A}\) （在 \(\mathbf{K}\) 上）由定理4是平展的。

关于平展代数的另一个刻画，见 V，第48页，命题1。

